\section{Weak coupling: what a proof of the mass gap still needs}\label{sec:weak}

The continuum limit of four-dimensional Yang--Mills theory is taken along the
tuning dictated by asymptotic freedom: for $\SU(N)$ the bare coupling
$g^2=2N/\beta_W$ tends to zero like $1/(2b_0\log(1/(a\Lambda)))$ with
$b_0=11N/(48\pi^2)$, the first coefficient of the $\beta$-function
\cite{GW73,Pol73} and a fixed physical scale $\Lambda$, so $\beta_W\to\infty$
logarithmically in $1/a$. Theorems~\ref{thm:A} and~\ref{thm:bpstrong} hold for
$\beta_W<N^2/29$ and $\beta_W<N^2/42$ respectively; at fixed $\beta$ the lattice
gap in physical units, $m_G(\beta)/a$, diverges as $a\to0$, and no continuum
limit exists at fixed $\beta$. The results of this paper therefore say nothing
about the continuum theory without statement (D). Three facts show what (D)
asks for.

\emph{(a) A cube condition is a mass.} By Theorem~\ref{thm:free} the best
constant for a free field is at most $s^2/(s^2+16d/9)$ with $s=m\ell$. At weak
coupling the leading approximation of the Wilson measure, in any gauge, is a
massless Gaussian field, for which the constant vanishes as $a\to0$ at fixed
$\ell$. Hence (D) does not follow from perturbation theory to any finite order:
it holds only if the interaction generates a mass of order $1/\ell$, and it is
a quantitative form of that mass generation.

\emph{(b) The cube must grow.} For $\SU(N)$, $\beta\ge1$ and the torus of side
$2n+4$, $\BP(n,\varpi)$ implies $\varpi\le C_Nn^4(1+\log\beta)/\beta$
\cite{BBLRS26}: on a small torus at weak coupling the Polyakov loops order, and
the cube dynamics can change their centre charge only through the curvature,
which is small. Along the tuning $n_k=\ell/a_k$ grows like a power of $1/a_k$
while $\beta_k$ grows like $\log(1/a_k)$, so this obstruction is consistent with
(D) but excludes every argument at fixed cube size.

\emph{(c) The group must be non-abelian.} For $G=\U(1)$ with the Villain
action at large $\beta$ the theory has a massless photon \cite{FS82}, and the
cube condition fails at every cube size, since by Theorem~\ref{thm:decay} it
would make the field-strength correlations summable \cite{BBLRS26}. Every proof
of (D) must therefore use the non-abelian structure of $G$, exactly as a proof
of the mass gap must.

\emph{The two ends.} Ba{\l}aban's renormalisation group
\cite{Bal84a,Bal85a,Bal89b} controls the effective actions of the lattice
theory from the scale $a$ up to scales at which the running coupling is still
small, uniformly in the lattice spacing, in finite volume. The theorems of
Secs.~\ref{sec:curvature}--\ref{sec:bpstrong} control the theory when the
coupling is large: they need $N\beta$ small, that is, $g$ large. The running
coupling grows towards the infrared, so a proof of (D) would follow the
renormalisation group from the weak-coupling regime, where Ba{\l}aban's
estimates hold, through the regime of coupling of order one, to a scale
$\ell$ at which the effective theory satisfies a strong-coupling criterion of
the kind proved here, and would then transport the resulting cube estimate
back to the fine lattice. Neither the middle step nor the final transport is
available: the effective actions are not of Wilson form, Ba{\l}aban's analysis
stops before the coupling is of order one, and no known criterion accepts the
effective theory at that scale. This gap between the two ends is the Yang--Mills
mass gap problem in the form isolated here.

Statement (E) is the construction problem. Ba{\l}aban's ultraviolet stability
bounds give compactness of suitably smeared observables in finite volume, but
the infinite-volume limit, Euclidean invariance and nontriviality of the limit
are open \cite{JW06}. Given (E), Theorem~\ref{thm:B} reduces the mass gap to
(D) with the explicit constant $\Delta\ge16\varpi/(625e^2\ell)$.

